\section{Discussion and limitations}
\label{sec:discussion}
\paragraph{What the evidence supports.}
On these 13 binary benchmarks, (i) a model-free change of design (difficulty stratification using source records) gives a consistent $\sim$17\% RMSE reduction; (ii) learned predictors give large further gains when the target is an interpolation of the source population, and become worse than random sampling---through a systematic downward bias---when the targets are stronger than every source model; (iii) per-target recalibration of an IRT predictor (Cal-IRT) or its pretest-shrunk version (ShrinkCal-IRT) retains much of the interpolation gain while staying better than random sampling in all our extrapolation settings, and Cal-IRT beats the AIPW estimator by 6--15\% in paired comparisons; (iv) the contribution of the variance-aware design is small ($\le$5--7\%, not always distinguishable from zero), and the uniform floor ($\lambda<1$) is safer; (v) selecting among methods by pseudo-shift validation inside the source set was not effective. In sample-size terms, and assuming the error of random sampling scales as $n^{-1/2}$, a relative RMSE of 0.70--0.76 corresponds to the accuracy of $\approx$1.7--2$\times$ as many random items. Practical reading: if one cannot rule out that a new model lies outside the span of the source models---the usual case for a new, stronger model---an uncorrected predictor can be worse than simply running random items, and a corrected or recalibrated estimator on a difficulty-stratified or VA-PPS coreset is the safer default; if the target is a sibling of known models, uncorrected predictors (or ShrinkCal-IRT) are better.

\paragraph{Limitations.}
\emph{Scope.} Only binary correctness on 13 benchmark matrices from two sources was studied; continuous or open-ended scores, chat/judge benchmarks and benchmarks with very few source models are not covered. The HELM benchmarks have only 79--83 models (about 20 targets per split), so their results are noisier. \emph{Shift is operationalised, not measured.} ``Stronger targets'' (top 25\% by accuracy) is one particular, fairly severe extrapolation; the ordering of methods under milder or different shifts (other capability profiles, different eras of models, new architectures) may differ, and the split is deterministic, so the intervals do not account for which models are strongest. Leaderboard models include many fine-tunes and merges of a few base models, which flatters interpolation (organisation hold-out only partly addresses this). \emph{Target accuracy is on the benchmark itself.} All targets are scored on the same finite item set; we say nothing about generalisation to unseen items, nor about evaluation noise in the underlying leaderboards (e.g.\ sampling temperature, prompt format). \emph{Baselines are re-implementations.} Anchor Points, tinyBenchmarks-style, D-optimal/metabench-style and AIPW baselines are our own simplified versions on binary data (e.g.\ $k$-means on binary vectors instead of confidence-based clustering, a fixed IRT dimension, no cross-validated mixing weights); they are not the original authors' code, and stronger tuned versions may close some gaps. \emph{Tuning.} IRT regularisation, the ability prior, the Platt penalty and $c_0$ were set on a separate development run using 9--11 of the benchmarks and fixed values; $c_0=0.5$ was not varied in the main run, and the VA-PPS $\lambda$ was fixed to the value that looked best for Cal-IRT in development, which turned out to be a worse choice than $\lambda=0.5$ for AIPW-type estimators. \emph{No guarantees.} Cal-IRT and ShrinkCal-IRT are heuristics: they are not unbiased (Cal-IRT has a small negative bias under extrapolation, $-0.5$ points), the shrinkage is a pretest rule without formal risk guarantees, and Proposition~\ref{prop:neyman} concerns Poisson sampling and a fixed predictor while we use systematic sampling with a proxy variance. \emph{Metric.} Relative RMSE averages errors over target models sharing a single coreset; ranking metrics are computed within a split and are not comparable across split types. \emph{Cost.} IRT fitting on the largest benchmark takes on the order of 10--30\,s per source set on one CPU core; per-target inference is negligible.

\paragraph{Broader impact.} Cheaper evaluation lowers costs, but over-trusting a predicted score (particularly for new, stronger models) can mislead model comparison; we recommend reporting the coreset size and an uncertainty interval, and not using predicted scores for decisions near ties.
